\subsubsection*{\texorpdfstring{Kramers \(\{-,-\}\) proper subgroups}
{Kramers \{-,-\} proper subgroups}}

The eight cases in this subsection obey
\(\Theta_1^2=\Theta_2^2=-I\) and
\(\Theta_1\Theta_2=\Theta_2\Theta_1\).
The two time reversals supplied by \(J_1\) and \(J_2\) therefore do not yet
form a Kramers pair.  An additional physical time reversal
\(\Theta_3=U_3\mc K\) selects generators that are invariant under a pair of
anticommuting transformations whose squares are both \(-I\).  This pair fixes
the sign of every Gaussian weight formed from the selected generators.

\medskip\noindent\textit{Bosons.}\quad
For the four bosonic cases, we use eight-component physical Nambu
representations with distinct \(J_1\) and \(J_2\).  The tensor factors are
ordered as Nambu, \({\rm f}_1\), and \({\rm f}_2\).  In each representation,
the maximal-subgroup equations give a 20-dimensional real vector space
spanned by Pauli tensors.  Every \(U_3\) displayed below gives
\(\Theta_3^2=-I\).  Imposing
    \(U_3^{-1}\overline{A_{\rm group}}U_3=A_{\rm group}\) retains the ten
    directions displayed below.

\medskip\noindent\textit{B11.}\quad
\(J_1^\trans gJ_1=g\), \(J_2^\trans gJ_2=g\), and
\(J_1J_2=J_2J_1\).
Choose
\begin{equation*}
\begin{aligned}
J_1&=I_2^{({\rm N})}\otimes I_2^{({\rm f}_1)}
       \otimes(\ii\sigma_y)^{({\rm f}_2)},\\
J_2&=I_2^{({\rm N})}\otimes\sigma_x^{({\rm f}_1)}
       \otimes(\ii\sigma_y)^{({\rm f}_2)},\\
U_3&=I_2^{({\rm N})}\otimes(\ii\sigma_y)^{({\rm f}_1)}
       \otimes I_2^{({\rm f}_2)} .
\end{aligned}
\end{equation*}
Here \(\Theta_1=J_1\mc K\), \(\Theta_2=J_2\mc K\), and
\(\Theta_3=U_3\mc K\).  Direct multiplication gives
\(\Theta_1\Theta_3=\Theta_3\Theta_1\) and
\(\Theta_2\Theta_3=-\Theta_3\Theta_2\).  Thus \(\Theta_2,\Theta_3\)
form the Kramers pair.  The 20-dimensional maximal-subgroup space is reduced to
\begin{align*}
A_{\rm group}\in\operatorname{span}_{\mathbb R}\{&
\ii I_2^{({\rm N})}\otimes I_2^{({\rm f}_1)}\otimes\sigma_y^{({\rm f}_2)},
\ \sigma_x^{({\rm N})}\otimes I_2^{({\rm f}_1)}\otimes I_2^{({\rm f}_2)},\\
&\ii\sigma_x^{({\rm N})}\otimes\sigma_x^{({\rm f}_1)}\otimes\sigma_x^{({\rm f}_2)},
\ \ii\sigma_x^{({\rm N})}\otimes\sigma_x^{({\rm f}_1)}\otimes\sigma_z^{({\rm f}_2)},\\
&\ii\sigma_y^{({\rm N})}\otimes I_2^{({\rm f}_1)}\otimes I_2^{({\rm f}_2)},
\ \sigma_y^{({\rm N})}\otimes\sigma_x^{({\rm f}_1)}\otimes\sigma_x^{({\rm f}_2)},\\
&\sigma_y^{({\rm N})}\otimes\sigma_x^{({\rm f}_1)}\otimes\sigma_z^{({\rm f}_2)},
\ \sigma_z^{({\rm N})}\otimes I_2^{({\rm f}_1)}\otimes I_2^{({\rm f}_2)},\\
&\ii\sigma_z^{({\rm N})}\otimes\sigma_x^{({\rm f}_1)}\otimes\sigma_x^{({\rm f}_2)},
\ \ii\sigma_z^{({\rm N})}\otimes\sigma_x^{({\rm f}_1)}\otimes\sigma_z^{({\rm f}_2)}
\}.
\end{align*}

\medskip\noindent\textit{B16.}\quad
\(J_1^\trans gJ_1=g\), \(J_2^\trans gJ_2=-g\), and
\(J_1J_2=-J_2J_1\).
Choose
\begin{equation*}
\begin{aligned}
J_1&=I_2^{({\rm N})}\otimes I_2^{({\rm f}_1)}
       \otimes(\ii\sigma_y)^{({\rm f}_2)},\\
J_2&=\sigma_z^{({\rm N})}\otimes(\ii\sigma_y)^{({\rm f}_1)}
       \otimes\sigma_x^{({\rm f}_2)},\\
U_3&=I_2^{({\rm N})}\otimes\sigma_x^{({\rm f}_1)}
       \otimes(\ii\sigma_y)^{({\rm f}_2)} .
\end{aligned}
\end{equation*}
The induced transformations are \(\Theta_1=J_1\mc K\),
\(\Theta_2=\ii J_2\mc K\), and \(\Theta_3=U_3\mc K\).
They obey \(\Theta_1\Theta_3=\Theta_3\Theta_1\) and
\(\Theta_2\Theta_3=-\Theta_3\Theta_2\).  The Kramers pair
\(\Theta_2,\Theta_3\) retains the following ten-dimensional space:
\begin{align*}
A_{\rm group}\in\operatorname{span}_{\mathbb R}\{&
\ii I_2^{({\rm N})}\otimes\sigma_x^{({\rm f}_1)}\otimes\sigma_y^{({\rm f}_2)},
\ \ii\sigma_x^{({\rm N})}\otimes I_2^{({\rm f}_1)}\otimes\sigma_x^{({\rm f}_2)},\\
&\sigma_x^{({\rm N})}\otimes\sigma_x^{({\rm f}_1)}\otimes I_2^{({\rm f}_2)},
\ \ii\sigma_x^{({\rm N})}\otimes\sigma_x^{({\rm f}_1)}\otimes\sigma_z^{({\rm f}_2)},\\
&\sigma_y^{({\rm N})}\otimes I_2^{({\rm f}_1)}\otimes\sigma_x^{({\rm f}_2)},
\ \ii\sigma_y^{({\rm N})}\otimes\sigma_x^{({\rm f}_1)}\otimes I_2^{({\rm f}_2)},\\
&\sigma_y^{({\rm N})}\otimes\sigma_x^{({\rm f}_1)}\otimes\sigma_z^{({\rm f}_2)},
\ \sigma_z^{({\rm N})}\otimes I_2^{({\rm f}_1)}\otimes I_2^{({\rm f}_2)},\\
&\ii\sigma_z^{({\rm N})}\otimes I_2^{({\rm f}_1)}\otimes\sigma_z^{({\rm f}_2)},
\ \ii\sigma_z^{({\rm N})}\otimes\sigma_x^{({\rm f}_1)}\otimes\sigma_x^{({\rm f}_2)}
\}.
\end{align*}

\begin{samepage}
\medskip\noindent\textit{B28.}\quad
\(J_1^\trans gJ_1=-g\), \(J_2^\trans gJ_2=g\), and
\(J_1J_2=-J_2J_1\).
Choose
\begin{equation*}
\begin{aligned}
J_1&=\sigma_z^{({\rm N})}\otimes I_2^{({\rm f}_1)}
       \otimes(\ii\sigma_y)^{({\rm f}_2)},\\
J_2&=I_2^{({\rm N})}\otimes(\ii\sigma_y)^{({\rm f}_1)}
       \otimes\sigma_x^{({\rm f}_2)},\\
U_3&=I_2^{({\rm N})}\otimes\sigma_x^{({\rm f}_1)}
       \otimes(\ii\sigma_y)^{({\rm f}_2)} .
\end{aligned}
\end{equation*}
The induced transformations are \(\Theta_1=\ii J_1\mc K\),
\(\Theta_2=J_2\mc K\), and \(\Theta_3=U_3\mc K\).
They obey \(\Theta_1\Theta_3=-\Theta_3\Theta_1\) and
\(\Theta_2\Theta_3=\Theta_3\Theta_2\).  The Kramers pair
\(\Theta_1,\Theta_3\) retains
\begin{align*}
A_{\rm group}\in\operatorname{span}_{\mathbb R}\{&
\ii I_2^{({\rm N})}\otimes\sigma_x^{({\rm f}_1)}\otimes\sigma_y^{({\rm f}_2)},
\ \ii\sigma_x^{({\rm N})}\otimes\sigma_y^{({\rm f}_1)}\otimes\sigma_y^{({\rm f}_2)},\\
&\ii\sigma_x^{({\rm N})}\otimes\sigma_z^{({\rm f}_1)}\otimes I_2^{({\rm f}_2)},
\ \sigma_x^{({\rm N})}\otimes\sigma_z^{({\rm f}_1)}\otimes\sigma_z^{({\rm f}_2)},\\
&\sigma_y^{({\rm N})}\otimes\sigma_y^{({\rm f}_1)}\otimes\sigma_y^{({\rm f}_2)},
\ \sigma_y^{({\rm N})}\otimes\sigma_z^{({\rm f}_1)}\otimes I_2^{({\rm f}_2)},\\
&\ii\sigma_y^{({\rm N})}\otimes\sigma_z^{({\rm f}_1)}\otimes\sigma_z^{({\rm f}_2)},
\ \sigma_z^{({\rm N})}\otimes I_2^{({\rm f}_1)}\otimes I_2^{({\rm f}_2)},\\
&\ii\sigma_z^{({\rm N})}\otimes I_2^{({\rm f}_1)}\otimes\sigma_z^{({\rm f}_2)},
\ \ii\sigma_z^{({\rm N})}\otimes\sigma_x^{({\rm f}_1)}\otimes\sigma_x^{({\rm f}_2)}
\}.
\end{align*}
\end{samepage}

\medskip\noindent\textit{B31.}\quad
\(J_1^\trans gJ_1=-g\), \(J_2^\trans gJ_2=-g\), and
\(J_1J_2=J_2J_1\).
Choose
\begin{equation*}
\begin{aligned}
J_1&=\sigma_z^{({\rm N})}\otimes I_2^{({\rm f}_1)}
       \otimes(\ii\sigma_y)^{({\rm f}_2)},\\
J_2&=\sigma_z^{({\rm N})}\otimes\sigma_x^{({\rm f}_1)}
       \otimes(\ii\sigma_y)^{({\rm f}_2)},\\
U_3&=I_2^{({\rm N})}\otimes(\ii\sigma_y)^{({\rm f}_1)}
       \otimes I_2^{({\rm f}_2)} .
\end{aligned}
\end{equation*}
The induced transformations are \(\Theta_1=\ii J_1\mc K\),
\(\Theta_2=\ii J_2\mc K\), and \(\Theta_3=U_3\mc K\).
They obey \(\Theta_1\Theta_3=-\Theta_3\Theta_1\) and
\(\Theta_2\Theta_3=\Theta_3\Theta_2\).  The Kramers pair
\(\Theta_1,\Theta_3\) retains
\begin{align*}
A_{\rm group}\in\operatorname{span}_{\mathbb R}\{&
\ii I_2^{({\rm N})}\otimes I_2^{({\rm f}_1)}\otimes\sigma_y^{({\rm f}_2)},
\ \sigma_x^{({\rm N})}\otimes I_2^{({\rm f}_1)}\otimes\sigma_x^{({\rm f}_2)},\\
&\sigma_x^{({\rm N})}\otimes I_2^{({\rm f}_1)}\otimes\sigma_z^{({\rm f}_2)},
\ \ii\sigma_x^{({\rm N})}\otimes\sigma_x^{({\rm f}_1)}\otimes I_2^{({\rm f}_2)},\\
&\ii\sigma_y^{({\rm N})}\otimes I_2^{({\rm f}_1)}\otimes\sigma_x^{({\rm f}_2)},
\ \ii\sigma_y^{({\rm N})}\otimes I_2^{({\rm f}_1)}\otimes\sigma_z^{({\rm f}_2)},\\
&\sigma_y^{({\rm N})}\otimes\sigma_x^{({\rm f}_1)}\otimes I_2^{({\rm f}_2)},
\ \sigma_z^{({\rm N})}\otimes I_2^{({\rm f}_1)}\otimes I_2^{({\rm f}_2)},\\
&\ii\sigma_z^{({\rm N})}\otimes\sigma_x^{({\rm f}_1)}\otimes\sigma_x^{({\rm f}_2)},
\ \ii\sigma_z^{({\rm N})}\otimes\sigma_x^{({\rm f}_1)}\otimes\sigma_z^{({\rm f}_2)}
\}.
\end{align*}

\medskip\noindent\textit{Fermions.}\quad
The four fermionic cases have four-component physical Nambu
representations.  The tensor factors are ordered as Nambu and flavor.  Each
maximal subgroup has a two-dimensional real generator space spanned by Pauli
tensors, and the additional time reversal retains one direction.

\begin{samepage}
\medskip\noindent\textit{F11.}\quad
\(J_1^\trans gJ_1=g\), \(J_2^\trans gJ_2=g\), and
\(J_1J_2=J_2J_1\).
Take
\begin{equation*}
J_1=I_2^{({\rm N})}\otimes(\ii\sigma_y)^{({\rm f})},
\qquad
J_2=\sigma_x^{({\rm N})}\otimes(\ii\sigma_y)^{({\rm f})},
\qquad
U_3=\ii(\ii\sigma_y)^{({\rm N})}\otimes I_2^{({\rm f})}.
\end{equation*}
Before imposing the \(U_3\) condition,
\(A_{\rm group}\in\operatorname{span}_{\mathbb R}
\{\ii I_2^{({\rm N})}\otimes\sigma_y^{({\rm f})},
\ii\sigma_x^{({\rm N})}\otimes\sigma_y^{({\rm f})}\}\).
It retains
\(A_{\rm group}=\kappa\,\ii I_2^{({\rm N})}\otimes\sigma_y^{({\rm f})}\),
with \(\kappa\in\mathbb R\).  For
\(\Theta_1=J_1\mc K\), \(\Theta_2=J_2\mc K\), and
\(\Theta_3=U_3\mc K\), direct multiplication gives
\(\Theta_1\Theta_3=-\Theta_3\Theta_1\) and
\(\Theta_2\Theta_3=\Theta_3\Theta_2\).
The retained generator is therefore invariant under the Kramers pair
\(\Theta_1,\Theta_3\).
\end{samepage}

\medskip\noindent\textit{F16.}\quad
\(J_1^\trans gJ_1=g\), \(J_2^\trans gJ_2=-g\), and
\(J_1J_2=-J_2J_1\).
Take
\begin{equation*}
J_1=I_2^{({\rm N})}\otimes(\ii\sigma_y)^{({\rm f})},
\qquad
J_2=(\ii\sigma_y)^{({\rm N})}\otimes\sigma_x^{({\rm f})},
\qquad
U_3=\sigma_x^{({\rm N})}\otimes(\ii\sigma_y)^{({\rm f})}.
\end{equation*}
Before imposing the \(U_3\) condition,
\(A_{\rm group}\in\operatorname{span}_{\mathbb R}
\{\ii\sigma_x^{({\rm N})}\otimes\sigma_y^{({\rm f})},
\ii\sigma_z^{({\rm N})}\otimes\sigma_x^{({\rm f})}\}\).
It retains
\(A_{\rm group}=\kappa\,\ii\sigma_x^{({\rm N})}
\otimes\sigma_y^{({\rm f})}\), with \(\kappa\in\mathbb R\).
For \(\Theta_1=J_1\mc K\), \(\Theta_2=\ii J_2\mc K\), and
\(\Theta_3=U_3\mc K\), one finds
\(\Theta_1\Theta_3=\Theta_3\Theta_1\) and
\(\Theta_2\Theta_3=-\Theta_3\Theta_2\).
The Kramers pair \(\Theta_2,\Theta_3\) leaves the retained generator invariant.

\medskip\noindent\textit{F28.}\quad
\(J_1^\trans gJ_1=-g\), \(J_2^\trans gJ_2=g\), and
\(J_1J_2=-J_2J_1\).
Take
\begin{equation*}
J_1=(\ii\sigma_y)^{({\rm N})}\otimes I_2^{({\rm f})},
\qquad
J_2=\sigma_x^{({\rm N})}\otimes(\ii\sigma_y)^{({\rm f})},
\qquad
U_3=I_2^{({\rm N})}\otimes(\ii\sigma_y)^{({\rm f})}.
\end{equation*}
Before imposing the \(U_3\) condition,
\(A_{\rm group}\in\operatorname{span}_{\mathbb R}
\{\ii I_2^{({\rm N})}\otimes\sigma_y^{({\rm f})},
\ii\sigma_z^{({\rm N})}\otimes I_2^{({\rm f})}\}\).
It retains
\(A_{\rm group}=\kappa\,\ii I_2^{({\rm N})}\otimes\sigma_y^{({\rm f})}\),
with \(\kappa\in\mathbb R\).  For
\(\Theta_1=\ii J_1\mc K\), \(\Theta_2=J_2\mc K\), and
\(\Theta_3=U_3\mc K\), direct multiplication gives
\(\Theta_1\Theta_3=-\Theta_3\Theta_1\) and
\(\Theta_2\Theta_3=\Theta_3\Theta_2\).
Thus \(\Theta_1,\Theta_3\) are the Kramers pair that preserves the retained
generator.

\medskip\noindent\textit{F31.}\quad
\(J_1^\trans gJ_1=-g\), \(J_2^\trans gJ_2=-g\), and
\(J_1J_2=J_2J_1\).
Take
\begin{equation*}
J_1=(\ii\sigma_y)^{({\rm N})}\otimes I_2^{({\rm f})},
\qquad
J_2=(\ii\sigma_y)^{({\rm N})}\otimes\sigma_x^{({\rm f})},
\qquad
U_3=I_2^{({\rm N})}\otimes(\ii\sigma_y)^{({\rm f})}.
\end{equation*}
Before imposing the \(U_3\) condition,
\(A_{\rm group}\in\operatorname{span}_{\mathbb R}
\{\ii\sigma_z^{({\rm N})}\otimes I_2^{({\rm f})},
\ii\sigma_z^{({\rm N})}\otimes\sigma_x^{({\rm f})}\}\).
It retains
\(A_{\rm group}=\kappa\,\ii\sigma_z^{({\rm N})}
\otimes\sigma_x^{({\rm f})}\), with \(\kappa\in\mathbb R\).
For \(\Theta_1=\ii J_1\mc K\), \(\Theta_2=\ii J_2\mc K\), and
\(\Theta_3=U_3\mc K\), one finds
\(\Theta_1\Theta_3=-\Theta_3\Theta_1\) and
\(\Theta_2\Theta_3=\Theta_3\Theta_2\).
The Kramers pair \(\Theta_1,\Theta_3\) preserves the retained generator.

The displayed spaces give Kramers \(\{-,-\}\) proper subgroups in all eight
cases.
